\begin{houseviewbox}[结论：业绩事件成立；正式 H1 前只做事件型观察]
雅化集团直接预告 2026H1 归母净利润为 11.00--13.00 亿元，其中 Q2 为 7.6118--9.6118 亿元，较 Q1 的 3.3882 亿元显著跳升。公司将改善归因于锂盐价格、销量与销售均价同步上行，以及矿--产--销协同、效率提升和成本控制。这个结论足以确认二季度利润弹性，却不足以确认 H2 的量、价、成本、资源供给、民爆并表或现金转换。

因此，当前行动不是用半年度预告推全年目标价，而是把正式 H1 报告定义为下一次决策点：若收入、现金流、应收、存货和分部说明能够解释利润，才启动全年模型；若利润成立而现金转换恶化，则预告只能被视为周期价格 $\beta$ 的阶段性收益。
\end{houseviewbox}

\begin{exhibitbox}[事件摘要与行动边界]
\begin{tabularx}{\exhibitboxwidth}{L{3.1cm}X L{3.0cm}X}
\toprule
项目 & 事实 / 结论 & 项目 & 行动含义 \\
\midrule
Q2 归母净利润 & 7.6118--9.6118 亿元，公司直接披露 & 业绩事件 & 已确认 Q2 弹性，不能全年化 \\
H1 归母净利润 & 11.00--13.00 亿元，未经审计预告 & 质量边界 & 等待正式收入、现金流及营运资本 \\
7 月 22 日收盘 & 16.79 元；市值 193.52 亿元 & 市场路径 & 相对 7 月 7 日 24.65 元回撤 31.9\% \\
外部 2026E EPS & 2.01--2.63 元，三份完整券商 PDF & 价格交叉核验 & 对应 6.38--8.35 倍；不是自有估值 \\
当前行动 & 事件型观察，正式 H1 前不建立基本面全年结论 & 有效期 & 至正式半年度报告披露后重开模型 \\
\bottomrule
\end{tabularx}
\end{exhibitbox}
\sourcenote{公司《2026 年第一季度报告》（2026-04-27）第 2 页；公司《2026 年半年度业绩预告》（2026-07-07）第 1 页；东财日线快照（2026-07-23 抓取）；国信、开源、东吴原始报告。}

\section{真正的争议：Q2 的弹性能否成为 H2 的持续性}

我们的判断是，行业价格/需求 $\beta$ 与公司经营因素共同驱动了 H1 改善，但前者证据更强。天齐、盛新和赣锋在其 2026Q1 披露中同样把锂价上行列为重要原因；雅化的 H1 预告还没有给出销量、实现售价、单位成本、自给率、套保和分部利润，不能量化公司特异性贡献。最强的反方是公司矿源、客户和民爆业务已经形成结构性盈利优势；要证明这一点，正式 H1 必须展示价格回落或波动时仍可持续的单位经济与现金转换，而不是只给合并归母净利润。

\begin{keyinsight}[本报告给投资者的可执行问题]
把问题从“Q2 利润是否高”改为“哪一个 H2 假设已经被验证”。国信的全年预测只要求 H2 单季平均归母约 5.08--6.08 亿元；东吴则要求约 8.66--9.66 亿元，接近或高于预告 Q2 区间。正式 H1 报告将先验证利润质量，无法直接验证 H2，但能决定这两类外部假设哪一类更有资格被保留。
\end{keyinsight}

\section{三项催化剂与三项风险}

\begin{tabularx}{\linewidth}{L{2.4cm}X L{2.4cm}X}
\toprule
正向催化剂 & 为什么重要 & 下行风险 & 为什么重要 \\
\midrule
H1 收入与归母匹配 & 证明利润并非只靠口径变化 & 现金流继续为负 & 将强化利润未转现的担忧 \\
锂盐量价成本披露 & 可拆分价格 $\beta$ 与公司经营贡献 & 应收/存货进一步占用 & 可能吞噬业绩修复的现金价值 \\
民爆分部收入/毛利验证 & 决定其是否真能缓冲周期 & 锂价回落、进口与库存压力 & 直接冲击 H2 持续性假设 \\
\bottomrule
\end{tabularx}
\sourcenote{公司 2025 年年报、2026Q1 报告、H1 业绩预告；SMM 2026H1 锂市场回顾；资料边界见附录。}
